\documentclass[10pt,a4paper]{amsart}
\usepackage{amsmath,amssymb,amsthm,graphicx}
% Retained minimal preamble for the Stacks Project chapter varieties.tex.
% Verbatim excerpt of the chapter preamble: Xy-pic setup (native preamble lines 14-18), the native theorem
% declarations with the mathematical macros (native preamble lines 154-217), and the project own no-op
% definitions of the reference/slogan/history display environments (native wrapper.tex lines 1-3) that the
% counted statement uses for its slogan. Copied mechanically; nothing added, removed or redefined.
% For commutative diagrams we use Xy-pic
\usepackage[all]{xy}

% We use 2cell for 2-commutative diagrams.
\xyoption{2cell}
% Theorem environments.
%
\theoremstyle{plain}
\newtheorem{theorem}[subsection]{Theorem}
\newtheorem{proposition}[subsection]{Proposition}
\newtheorem{lemma}[subsection]{Lemma}

\theoremstyle{definition}
\newtheorem{definition}[subsection]{Definition}
\newtheorem{example}[subsection]{Example}
\newtheorem{exercise}[subsection]{Exercise}
\newtheorem{situation}[subsection]{Situation}

\theoremstyle{remark}
\newtheorem{remark}[subsection]{Remark}
\newtheorem{remarks}[subsection]{Remarks}

\numberwithin{equation}{subsection}

% Macros
%
\def\lim{\mathop{\mathrm{lim}}\nolimits}
\def\colim{\mathop{\mathrm{colim}}\nolimits}
\def\Spec{\mathop{\mathrm{Spec}}}
\def\Hom{\mathop{\mathrm{Hom}}\nolimits}
\def\Ext{\mathop{\mathrm{Ext}}\nolimits}
\def\SheafHom{\mathop{\mathcal{H}\!\mathit{om}}\nolimits}
\def\SheafExt{\mathop{\mathcal{E}\!\mathit{xt}}\nolimits}
\def\Sch{\mathit{Sch}}
\def\Mor{\mathop{\mathrm{Mor}}\nolimits}
\def\Ob{\mathop{\mathrm{Ob}}\nolimits}
\def\Sh{\mathop{\mathit{Sh}}\nolimits}
\def\NL{\mathop{N\!L}\nolimits}
\def\CH{\mathop{\mathrm{CH}}\nolimits}
\def\proetale{{pro\text{-}\acute{e}tale}}
\def\etale{{\acute{e}tale}}
\def\QCoh{\mathit{QCoh}}
\def\Ker{\mathop{\mathrm{Ker}}}
\def\Im{\mathop{\mathrm{Im}}}
\def\Coker{\mathop{\mathrm{Coker}}}
\def\Coim{\mathop{\mathrm{Coim}}}

% Boxtimes
%
\DeclareMathSymbol{\boxtimes}{\mathbin}{AMSa}{"02}

%
% Macros for moduli stacks/spaces
%
\def\QCohstack{\mathcal{QC}\!\mathit{oh}}
\def\Cohstack{\mathcal{C}\!\mathit{oh}}
\def\Spacesstack{\mathcal{S}\!\mathit{paces}}
\def\Quotfunctor{\mathrm{Quot}}
\def\Hilbfunctor{\mathrm{Hilb}}
\def\Curvesstack{\mathcal{C}\!\mathit{urves}}
\def\Polarizedstack{\mathcal{P}\!\mathit{olarized}}
\def\Complexesstack{\mathcal{C}\!\mathit{omplexes}}
% \Pic is the operator that assigns to X its picard group, usage \Pic(X)
% \Picardstack_{X/B} denotes the Picard stack of X over B
% \Picardfunctor_{X/B} denotes the Picard functor of X over B
\def\Pic{\mathop{\mathrm{Pic}}\nolimits}
\def\Picardstack{\mathcal{P}\!\mathit{ic}}
\def\Picardfunctor{\mathrm{Pic}}
\def\Deformationcategory{\mathcal{D}\!\mathit{ef}}
\newenvironment{reference}{}{}
\newenvironment{slogan}{}{}
\newenvironment{history}{}{}
\usepackage{cleveref}
\crefname{originaltheorem}{original theorem}{original theorems}
\Crefname{originaltheorem}{Original theorem}{Original theorems}
\makeatletter
\newlabel{algebra-theorem-nullstellensatz}{{00FV}{source}{}{sourcepointer-1}{}}
\newlabel{algebra-theorem-nullstellensatz@cref}{{[originaltheorem][0][]00FV}{[1][0][]source}{}{}{}}
\makeatother
\crefname{originallemma}{original lemma}{original lemmas}
\Crefname{originallemma}{Original lemma}{Original lemmas}
\makeatletter
\newlabel{morphisms-lemma-Jacobson-universally-Jacobson}{{02J5}{source}{}{sourcepointer-2}{}}
\newlabel{morphisms-lemma-Jacobson-universally-Jacobson@cref}{{[originallemma][0][]02J5}{[1][0][]source}{}{}{}}
\makeatother
\crefname{originallemma}{original lemma}{original lemmas}
\Crefname{originallemma}{Original lemma}{Original lemmas}
\makeatletter
\newlabel{morphisms-lemma-jacobson-finite-type-points}{{01TB}{source}{}{sourcepointer-3}{}}
\newlabel{morphisms-lemma-jacobson-finite-type-points@cref}{{[originallemma][0][]01TB}{[1][0][]source}{}{}{}}
\makeatother
\crefname{originallemma}{original lemma}{original lemmas}
\Crefname{originallemma}{Original lemma}{Original lemmas}
\makeatletter
\newlabel{morphisms-lemma-ubiquity-Jacobson-schemes}{{02J6}{source}{}{sourcepointer-4}{}}
\newlabel{morphisms-lemma-ubiquity-Jacobson-schemes@cref}{{[originallemma][0][]02J6}{[1][0][]source}{}{}{}}
\makeatother
\crefname{originaldefinition}{original definition}{original definitions}
\Crefname{originaldefinition}{Original definition}{Original definitions}
\makeatletter
\newlabel{properties-definition-jacobson}{{01P2}{source}{}{sourcepointer-5}{}}
\newlabel{properties-definition-jacobson@cref}{{[originaldefinition][0][]01P2}{[1][0][]source}{}{}{}}
\makeatother
\crefname{originallemma}{original lemma}{original lemmas}
\Crefname{originallemma}{Original lemma}{Original lemmas}
\makeatletter
\newlabel{properties-lemma-locally-jacobson}{{01P4}{source}{}{sourcepointer-6}{}}
\newlabel{properties-lemma-locally-jacobson@cref}{{[originallemma][0][]01P4}{[1][0][]source}{}{}{}}
\makeatother
\title{A scheme locally of finite type over a field is Jacobson with algebraic residue fields at closed points}
\author{The Stacks Project Authors}
\date{stacks-project:varieties.tex:local-snapshot-20261011; GNU FDL 1.2 or later}
\begin{document}
\maketitle
\noindent\textit{Editorial note.} All delivered mathematical text is the author's own, copied from the pinned Stacks Project chapter file varieties.tex; no statement or proof step was rewritten, reordered or removed.\par\medskip
\noindent\textit{Editorial note.} The author gives two independent routes to the result, the second by reference to the morphisms chapter; both sentences are reproduced verbatim and the cited results are carried as navigation stubs.\par\medskip
\noindent\textit{Editorial note.} The counted spans carry 6 navigation stubs and no bibliography citation.\par\medskip
\noindent\textit{Editorial note.} The preamble is the chapter's retained minimal preamble (sources/stacks-varieties/retained-preamble.tex lines 1-77), a verbatim excerpt of the native preamble.\par\medskip
\noindent\textit{Editorial note.} The label lemma-locally-finite-type-Jacobson is defined only in varieties.tex, so no cross-chapter restatement had to be weighed for this card.\par\medskip
\noindent\textit{Allowed standard prerequisites (editorial).}\begin{itemize}
\item Jacobson property is local on an affine open covering. Properties, Lemma properties-lemma-locally-jacobson and Definition properties-definition-jacobson, invoked at proof lines 2682-2684 and named in the statement. Supplies the reduction to the affine case for both properties.
\item Hilbert Nullstellensatz. Algebra, Theorem algebra-theorem-nullstellensatz, invoked at proof lines 2686-2688. Supplies the affine case of both assertions.
\item Jacobson schemes and finite type points in the morphisms chapter. Morphisms, Lemmas morphisms-lemma-jacobson-finite-type-points, morphisms-lemma-Jacobson-universally-Jacobson and morphisms-lemma-ubiquity-Jacobson-schemes, invoked at proof lines 2689-2692. Supplies the author's second, independent route to the same conclusion.
\end{itemize}\medskip
\ifdefined\hypertarget\hypertarget{sourcepointer-1}{}\fi
\noindent\textit{Source reference.} Original theorem 00FV: sources/stacks-chow/tags.txt, source lines 585--585. Pointer outside the counted span, invoked once for context; the load-bearing argument of the counted statement is the author's proof reproduced in full in the proof span.\par\medskip
\ifdefined\hypertarget\hypertarget{sourcepointer-2}{}\fi
\noindent\textit{Source reference.} Original lemma 02J5: sources/stacks-chow/tags.txt, source lines 3324--3324. Pointer outside the counted span, invoked once for context; the load-bearing argument of the counted statement is the author's proof reproduced in full in the proof span.\par\medskip
\ifdefined\hypertarget\hypertarget{sourcepointer-3}{}\fi
\noindent\textit{Source reference.} Original lemma 01TB: sources/stacks-chow/tags.txt, source lines 2368--2368. Pointer outside the counted span, invoked once for context; the load-bearing argument of the counted statement is the author's proof reproduced in full in the proof span.\par\medskip
\ifdefined\hypertarget\hypertarget{sourcepointer-4}{}\fi
\noindent\textit{Source reference.} Original lemma 02J6: sources/stacks-chow/tags.txt, source lines 3325--3325. Pointer outside the counted span, invoked once for context; the load-bearing argument of the counted statement is the author's proof reproduced in full in the proof span.\par\medskip
\ifdefined\hypertarget\hypertarget{sourcepointer-5}{}\fi
\noindent\textit{Source reference.} Original definition 01P2: sources/stacks-chow/tags.txt, source lines 2215--2215. Pointer outside the counted span, invoked once for context; the load-bearing argument of the counted statement is the author's proof reproduced in full in the proof span.\par\medskip
\ifdefined\hypertarget\hypertarget{sourcepointer-6}{}\fi
\noindent\textit{Source reference.} Original lemma 01P4: sources/stacks-chow/tags.txt, source lines 2217--2217. Pointer outside the counted span, invoked once for context; the load-bearing argument of the counted statement is the author's proof reproduced in full in the proof span.\par\medskip
\begin{lemma}
\label{lemma-locally-finite-type-Jacobson}
Let $X$ be a scheme which is locally of finite type over $k$.
Then
\begin{enumerate}
\item for any closed point $x \in X$ the extension $\kappa(x)/k$
is algebraic, and
\item $X$ is a Jacobson scheme
(Properties, Definition \ref{properties-definition-jacobson}).
\end{enumerate}
\end{lemma}

\begin{proof}
A scheme is Jacobson if and only if it has an affine open covering
by Jacobson schemes, see
Properties, Lemma \ref{properties-lemma-locally-jacobson}.
The property on residue fields at closed points is also local on $X$.
Hence we may assume that $X$ is affine. In this case the result
is a consequence of the Hilbert Nullstellensatz, see
Algebra, Theorem \ref{algebra-theorem-nullstellensatz}.
It also follows from a combination of
Morphisms, Lemmas \ref{morphisms-lemma-jacobson-finite-type-points},
\ref{morphisms-lemma-Jacobson-universally-Jacobson}, and
\ref{morphisms-lemma-ubiquity-Jacobson-schemes}.
\end{proof}

\end{document}
